
Il s'%
ensuit (V) que @@M0@@ est la @@M1@@-ème fonction symétrique élémentaire en les @@M2@@ (@@M3@@), donc est un polynôme homogène de degré @@M4@@ à coefficients entiers dans les variables @@M5@@ (@@M6@@). Appliquant à nouveau l'hypothèse de récurrence et l'associativité de la loi (1.3.1) on est ainsi ramené à démontrer que, pour tout @@M7@@-module inversible @@M8@@, on a
@@M9@@
Le schéma @@M10@@ étant intègre, @@M11@@ et @@M12@@ sont isomorphes à des sous-modules du @@M13@@-module constant @@M14@@ des fonctions rationnelles sur @@M15@@. Le @@M16@@-module somme @@M17@@ est cohérent, et on trouve
\begin{align*}
(\cl^{\bullet}(\underline{L})-1)\cl_\bullet(\underline{O}_X)
&=\cl_\bullet(\underline{L})-\cl_\bullet(\underline{O}_X)\\
&=(\cl_\bullet(\underline{F})-\cl_\bullet(\underline{O}_X))-(\cl_\bullet(\underline{F})-\cl_\bullet(\underline{L}))\\
&=\cl_\bullet(\underline{F}/\underline{O}_X)-\cl_\bullet(\underline{F}/\underline{L}).
\end{align*}
Mais les fibres au point générique des faisceaux @@M18@@ et @@M19@@ sont nulles, donc leurs supports dans @@M20@@ sont rares (@@M21@@ étant irréductible), donc de dimension @@M22@@. L'assertion (1.3.6) en découle immédiatement.

II. Reste à démontrer le cas général. On introduit le fibré en drapeaux
@@M23@@
de @@M24@@ sur @@M25@@, où on peut trouver des @@M26@@-modules inversibles @@M27@@ tels que
@@M28@@
Soient @@M29@@ le point générique de @@M30@@, @@M31@@ le corps résiduel de @@M32@@. Alors on peut trouver un point rationnel @@M33@@ de la fibre générique @@M34@@ de @@M35@@ 
